\documentclass[../main.tex]{subfiles}
\begin{document}

\section{\textit{Screening} termodinámico}

% ===============================================
% DESCRIPCIÓN DEL MODELO
% ===============================================

\subsection{Descripción del modelo}
El primer paso de este TFG es realizar un \textit{screening} termodinámico para identificar las posibles mezclas que superen al propano (R290) en rendimiento energético.

El \textit{screening} consistirá en la combinación de los posibles refrigerantes de trabajo en mezclas de tres en tres. Con cada mezcla se calcularán las propiedades termodinámicas del ciclo de compresión con incrementos y decrementos en las proporciones de los refrigerantes de un 5\%, resultando en un total de 231 posibles combinaciones para cada tres refrigerantes. Posteriormente, se filtrarán los resultados en función de unas restricciones (ver \cref{subsec:restricciones} donde se muestran las restricciones) y de estos resultados se calcularán la mejora del COP y la variación de VHC.

Si la mezcla resulta mejor que el propano en cuento a eficiencia y está dentro del rango de variación de VHC, entonces, se reevaluará con incrementos de menor proporción (1\%) y de esta reevaluación se elegirá la proporción que maximice el COP y esté también dentro del rango de variación de la VHC. Véase la \autoref{fig:diagrama_screening} donde se representa el diagrama de flujo del \textit{scrrening}.

El comportamiento de cada combinación se simulará considerando su operación en un ciclo básico de vapor. Este involucrará una compresión con un rendimiento fijo, una transformación isobárica en cada intercambiador de calor y una expansión isoentálpica en la válvula de expansión. Cada combinación se simulará en cuatro escenarios diferentes, dos vienen tpificados por la normativa UNE-EN 14511-3 \cite{une14511} para ensayos de bombas de calor agua-agua y otros dos tienen condiciones parecidas a la normativa pero se han añadido para reducir la incertidumbre de medida en el evaporador.

\newpage

\begin{figure}[H]
    \centering
    \includegraphics[width=0.85\textwidth]{imagenes/diagrama_screening.pdf}
    \caption{Diagrama del \textit{screening} termodinámico}
    \label{fig:diagrama_screening}
\end{figure}

% ===============================================
% CONDICIONES DE SIMULACIÓN
% ===============================================

\subsection{Condiciones de simulación}

\subsubsection{Fluidos}

Tras considerar una gran serie fluidos naturales junto al grupo GIT, se ha considerado realizar el \textit{screening} termodinámico de los refrigerantes de la \autoref{tab:fluidos_screening} debido a su dispnibilidad tanto económica como técnica ya que son bien conocidos y sus características en la base de datos REFPROP es precisa.

\begin{table}[ht]
    \centering
    \begin{tabular}{c|c|c|c|c|c|c|}
    \cline{2-7}
    & \textbf{Familia} & \textbf{\begin{tabular}[c]{@{}c@{}}Código\\ ASHRAE\end{tabular}} & \textbf{GWP} & \textbf{ODP} & \textbf{NBP (ºC)} & \textbf{\begin{tabular}[c]{@{}c@{}}Código de\\ peligrosidad\end{tabular}} \\ \hline
    \multicolumn{1}{|c|}{\textbf{Propano}} & HC & R290 & 3.3 - 9.5 & 0 & -42.1 & A3 \\ \hline
    \multicolumn{1}{|c|}{\textbf{Butano}} & HC & R600 & 4.0 - 6.5 & 0 & -0.49 & A3 \\ \hline
    \multicolumn{1}{|c|}{\textbf{Isobutano}} & HC & R600a & 3 & 0 & -11.75 & A3 \\ \hline
    \multicolumn{1}{|c|}{\textbf{CO2}} & Óxido & R744 & 1 & 0 & - & A1 \\ \hline
    \multicolumn{1}{|c|}{\textbf{Dimetil Éter}} & HC & RE170 & 1 & 0 & -24.78 & A3 \\ \hline
    \multicolumn{1}{|c|}{\textbf{Etano}} & HC & R170 & 5.5 - 10.2 & 0 & -88.58 & A3 \\ \hline
    \multicolumn{1}{|c|}{\textbf{Etileno}} & HO & R1150 & 3.7 & 0 & -103.77 & A3 \\ \hline
    \multicolumn{1}{|c|}{\textbf{Propileno}} & HO & R1270 & 1.8 & 0 & -57.62 & A3 \\ \hline
    \multicolumn{1}{|c|}{\textbf{Metano}} & HC & R50 & 28 & 0 & -161.48 & A3 \\ \hline
    \end{tabular}
    \caption{Tabla de fluidos que se estudiarán en el \textit{screening}}
    \label{tab:fluidos_screening}
\end{table}

\subsubsection{Variables}
Para el cálculo termodinámico del ciclo las principales variables tomarán los siguientes valores para cada uno de los escenarios:

En la \autoref{fig:condiciones_constantes} se muestran las variables que son constantes para todos los ensayos y en la \autoref{fig:cond_une_teorico} se muestran las condiciones que varian en función de los 4 diferentes escenarios que se analizan.

\begin{table}[H]
\centering
\begin{tabular}{|c|c|c|c|c|}
\hline
\textbf{SUB} & \textbf{SH} & $\textbf{Ap}_\textbf{k}$ & $\textbf{Ap}_\textbf{o}$  & $\mathbf{\eta}_\textbf{comp}$\\ \hline
1 ºC & 6 ºC& 6.5 ºC& 3 ºC&0.6\\ \hline
\end{tabular}
\caption{Condiciones constantes en los ensayos}
\label{fig:condiciones_constantes}
\end{table}

\begin{table}[H]
    \centering
    \begin{tabular}{c|cc|cc|}
\cline{2-5}
 & \multicolumn{2}{c|}{\textbf{Condensador}} & \multicolumn{2}{c|}{\textbf{Evaporador}} \\ \hline
\multicolumn{1}{|c|}{\textbf{Tipo de ensayo}} & \multicolumn{1}{c|}{\textbf{Entrada}} & \textbf{Salida} & \multicolumn{1}{c|}{\textbf{Entrada}} & \textbf{Salida} \\ \hline
\multicolumn{1}{|c|}{Temperatura alta} & \multicolumn{1}{c|}{55 ºC} & 65 ºC & \multicolumn{1}{c|}{0 ºC} & -3 ºC \\ \hline
\multicolumn{1}{|c|}{Temperatura media} & \multicolumn{1}{c|}{47 ºC} & 55 ºC & \multicolumn{1}{c|}{0 ºC} & -3 ºC \\ \hline
\end{tabular}
    \caption{Temperaturas de los ensayos}
    \label{fig:cond_une_teorico}
\end{table}

% ===============================================
% ECUACIONES DE CÁLCULO
% ===============================================

\subsection{Ecuaciones de cálculo}
La presión de condensación se ha encontrado usando tanto la temperatura de entrada del fluido como el approach de temperatura y el subcooling:
\begin{equation}
p_{k} = f(t = t_{hw\_in} + Ap_{k} + SUB,\ x_{v} = 0)
\end{equation}
Una vez conocido el punto termodinámico de la salida del condensador, se puede hallar de forma similar el punto posterior a la válvula de expansión:
\begin{equation}
p_{O} = f(t = t_{cw\_out} - Ap_{O},\ h = h_{k,out})
\end{equation}
A partir de estas presiones se pueden hallar un par de magnitudes que definan termodinámicamente los puntos de salida del condensador y de entrada del evaporador.
\begin{gather}
\mathrm{h}_{k,out} = f(p = p_{k},\ t = t_{hw\_in} + Ap_{k})
\label{eq:punto_k,out}
\\
\mathrm{h}_{O,in} = f(p = p_{O},\ h = h_{k,out})
\label{eq:punto_O,in}
\end{gather}
Obteniendo la temperatura de vapor saturado en el evaporador, se puede hallar el punto de salida del evaporador:
\begin{gather}
t_{vap,sat} = f(p = p_{O},\ x_{v} = 1)
\label{eq:t_vap_sat}
\\
\mathrm{h}_{O,out} = f(p = p_{O},\ t = t_{vap,sat})
\label{eq:punto_O,out}
\end{gather}
Conociendo el rendimiento del compresor, se puede hallar el punto a la salida de este, que es el mismo que el de la entrada del condensador:
\begin{gather}
h_{c,out,s} = f(p = p_{k},\ s = s_{c,in}),\quad
\label{eq:h_c,out,s}
\\
h_{c,out} = h_{c,in} + \frac{h_{c,out,s} - h_{c,in}}{\eta_{comp}}
\label{eq:h_c,out}
\\
\mathrm{h}_{k,in} = f(p = p_{k},\ h = h_{c,out})
\label{eq:punto_k,in}
\end{gather}

A continuación, se muestran las expresiones que se usan para calcular los parámetros del ciclo termodinámico:
\begin{gather}
Glide_{O}=t(p = p_{O},\ x_{v} = 1) - t(p = p_{O},\ x_{v} = x_{v,O,in})
\label{eq:glide_O}
\\
Glide_{k}=t(p = p_{k},\ x_{v} = 1) - t(p = p_{k},\ x_{v} = 0)
\label{eq:glide_k}
\\
VHC = \frac{h_{k,in} - h_{k,ou}}{\nu_{suc}}
\label{eq:VHC}
\\
COP = \frac{h_{O,out}-h_{O,in}}{h_{k,in}-h_{k,out}}
\label{eq:COP}
\\
\frac{\dot{m}_{hw}}{\dot{m}_{ref}}=\frac{h_{O,in}-h_{O,out}}
{c_{p}\cdot \left( t_{hw,out}-t_{hw,in} \right)}
\label{eq:caudales}
\\
t_{hw,pinch} = t_{hw,out} - \frac{\dot{m}_{ref}}{\dot{m}_{hw}\cdot c_{p}} \cdot 
\left( h_{k,in} - h(p = p_{k},\ x_{v} = 1)\right)
\label{eq:t_hw_pinch}
\\
t_{pinch} = t(p = p_{k},\ x_{v} = 1) - t_{hw,pinch}
\label{t_pinch}
\end{gather}

El software usado para el cálculo de las propiedades termodinámicas ha sido REFPROP v.10.0 \cite{refprop}

% ===============================================
% RESTRICCIONES
% ===============================================

\subsection{Restricciones}
\label{subsec:restricciones}
El proceso incluye filtros que excluyen las mezclas que no cumplen los siguientes requisitos:
\begin{itemize}
    \item Glide de condensación ($Glide_{k}$) y glide efectivo de evaporación 
    ($Glide_{O}$) menor a $10\,^{\circ}\mathrm{C}$, ya que la transferencia de 
    calor se reduce al pasar estos umbrales.
    \item Temperatura de descarga del compresor menor a $130\,^{\circ}\mathrm{C}$,
    ya que este no está pensado para soportar temperaturas tan elevadas.
    \item Presión de descarga del compresor menor a $28\ \mathrm{bar}$, ya que esta es
    su presión máxima de operación.
    \item Valores de $VHC$ que estén en el intervalo $\pm 30\%\ VHC_{R290}$,  
    para ser compatibles con los compresores de propano y no afecar significativamente a su capacidad calorífica.
    \item Presión de descarga del compresor menor que la presión transcrítica del fluido, ya que la instalación no está preparada para este estado.
\end{itemize}
Además, si la diferencia de temperatura entre el fluido y el agua en el \textit{pinch point} es menor a $1\,^{\circ}\mathrm{C}$, se volverá a calcular el ciclo pero aumentando en $0.5\,^{\circ}\mathrm{C}$ el valor del approach de temperatura del condensador hasta que el ciclo cumpla este criterio.

Posteriormente, todas las mezclas que superen los filtros se reevaluarán pero ahora con una variación del 1\% en su composición para encontrar la composición que maximiza el $COP$.

% ===============================================
% RESULTADOS TEÓRICOS
% ===============================================
\newpage
\subsection{Resultados teóricos}

\begin{table}[ht]
    \centering
    
\begin{tabular}{|ccc|ccc|c|c|c|c|}
\hline
\multicolumn{3}{|c|}{\multirow{2}{*}{\textbf{Fluidos}}} & \multicolumn{3}{c|}{\multirow{2}{*}{\textbf{Prporciones}}} & \multicolumn{1}{c|}{\multirow{2}{*}{\textbf{COP}}} & \multicolumn{1}{c|}{\multirow{2}{*}{\textbf{VHC}}} & \multicolumn{2}{c|}{\textbf{Glide (ºC)}} \\ \cline{9-10} 
\multicolumn{3}{|c|}{} & \multicolumn{3}{c|}{} & \multicolumn{1}{c|}{} & \multicolumn{1}{c|}{} & \multicolumn{1}{c|}{\textbf{cond}} & \multicolumn{1}{c|}{\textbf{evap}} \\ \hline
R1270 & RE170 & R600 & 4\% & 72\% & 24\% & +16.4\% & -27.1\% & 3.0 & 6.1 \\ \hline
R290 & RE170 & R600 & 15\% & 52\% & 33\% & +13.8\% & -28.0\% & 6.4 & 8.4 \\ \hline
R1270 & RE170 & R600a & 19\% & 65\% & 16\% & +13.6\% & -14.2\% & 2.1 & 4.5 \\ \hline
RE170 & R744 & - & 99\% & 1\% & - & +10.9\% & -26.0\% & 1.2 & 0.3 \\ \hline
RE170 & R1270 & - & 93\% & 7\% & - & +10.9\% & -23.1\% & 1.1 & 0.8 \\ \hline
RE170 & R290 & - & 95\% & 5\% & - & +10.8\% & -24.1\% & 1.0 & 0.7 \\ \hline
RE170 & R170 & - & 98\% & 2\% & - & +10.8\% & -25.1\% & 1.4 & 0.5 \\ \hline
RE170 & R1150 & - & 99\% & 1\% & - & +10.6\% & -25.9\% & 1.5 & 0.4 \\ \hline
RE170 & R600 & - & 99\% & 1\% & - & +10.2\% & -25.1\% & 0.1 & 0 \\ \hline
RE170 & - & - & 100\% & - & - & +9.7\% & -27.8\% & 0 & 0 \\ \hline
\end{tabular}

    \caption{Resultado teórico con temperaturas altas (ensayo tipo 1)}
    \label{tab:res_teo_alta}
\end{table}

\begin{table}[ht]
    \centering
    
\begin{tabular}{|ccc|ccc|c|c|c|c|}
\hline
\multicolumn{3}{|c|}{\textbf{Fluidos}} & \multicolumn{3}{c|}{\textbf{Prporciones}} & \textbf{$\Delta$ COP} & \textbf{$\Delta$ VHC} & \textbf{glide\_k (ºC)} & \textbf{glide\_o (ºC)} \\ \hline
\multicolumn{1}{|c|}{R1270} & \multicolumn{1}{c|}{RE170} & R600 & \multicolumn{1}{c|}{3\%} & \multicolumn{1}{c|}{66\%} & 31\% & 14.3\% & -31\% & 3.9 & 7.4 \\ \hline
\multicolumn{1}{|c|}{R290} & \multicolumn{1}{c|}{RE170} & R600 & \multicolumn{1}{c|}{15\%} & \multicolumn{1}{c|}{51\%} & 34\% & 11.9\% & -29\% & 6.9 & 9.3 \\ \hline
\multicolumn{1}{|c|}{R1270} & \multicolumn{1}{c|}{RE170} & R600a & \multicolumn{1}{c|}{19\%} & \multicolumn{1}{c|}{65\%} & 16\% & 11.4\% & -15\% & 2.1 & 4.9 \\ \hline
\multicolumn{3}{|c|}{RE170} & \multicolumn{3}{c|}{100\%} & 7.9\% & -30\% & 0 & 0 \\ \hline
\multicolumn{3}{|c|}{\begin{tabular}{c|c} R290 & R600 \end{tabular}} & \multicolumn{3}{c|}{\begin{tabular}{c|c} 70\% & 30\% \end{tabular}} & 3.1\% & -26\% & 10 & 10 \\ \hline
\multicolumn{1}{|c|}{R1270} & \multicolumn{1}{c|}{R600} & R600a & \multicolumn{1}{c|}{60\%} & \multicolumn{1}{c|}{20\%} & 20\% & 2.7\% & -25\% & 10 & 8.7 \\ \hline
\multicolumn{1}{|c|}{} & \multicolumn{1}{c|}{R1270} & R600a & \multicolumn{1}{c|}{} & \multicolumn{1}{c|}{50\%} & 50\% & 2.3\% & -30\% & 6.4 & 4.8 \\ \hline
\end{tabular}

    \caption{Resultado teórico con temperaturas medias (ensayo tipo 4)}
    \label{tab:res_teo_media}
\end{table}


\FloatBarrier

En las tablas superiores (\autoref{tab:res_teo_alta} y \autoref{tab:res_teo_media}) se muestran las diez mejores mezclas de cada condición de simulación. Cabe destacar que no se están representando los resultados de los ensayos tipo 2 y tipo 3 ya que al solamente variar las temperaturas del fluido secundario en el evaporador no representan un cambio en los resultados debido a las limitaciones del modelo teórico.

% ===============================================
% ANÁLISIS DE RESULTADOS TEÓRICOS
% ===============================================

\subsection{Análisis de resultados}

Como se puede observar en la \autoref{fig:teo_COP_VHC_alta} y en la \autoref{fig:teo_COP_VHC_media}, la gran mayoría de los fluidos con mayor incremento de COP tienen una merma considerable de capacidad calorífica volumétrica, esto no es deseable, ya que implica un aumento en la frecuencia de trabajo del compresor, que puede tener como consecuencia una posible reducción de rendimiento del compresor, ya que se está variando el punto de trabajo.

Cabe destacar que se ha añadido una mezcla a los resultados (R290 / RE170, 85\% / 15\%), esta pertenece al estudio de Mickoleit et al. \cite{Mickoleit_2026} donde se realiza otro \textit{screening} termodinámico cuyo mejor resultado fue esta mezcla. Se ha decidido incluir en este trabajo ya que el estudio propuesto por Mickoleit et al. no contiene datos exerimentales y estudiarla experimentalmente podría \hl{NO SE QUE POSAR, REFORMULAR}. Además, como mencionan en su artículo, esta mezcla puede ser interesante, ya que presenta propiedades similares a las que tendría un fluido puro sin serlo.

\begin{figure}[ht]
    \centering
    \includegraphics[width=\textwidth]{imagenes/t_alta_COP_VHC.png}
    \caption{Comparativa de COP y VHC de los resultados teóricos (temperatura alta)}
    \label{fig:teo_COP_VHC_alta}
\end{figure}

\begin{figure}[ht]
    \centering
    \includegraphics[width=\textwidth]{imagenes/t_media_COP_VHC.png}
    \caption{Comparativa de COP y VHC de los resultados teóricos (temperatura media)}
    \label{fig:teo_COP_VHC_media}
\end{figure}

\FloatBarrier

\subsubsection{Temperatura alta}

Los resultados en esta configuración están dominados por tres mezclas ternarias seguidos de de mezclas que mezclas que mayoritariamente contienen dimetil con una pequeña proporción de otro fluido. Por último, se encuentra la mezcla de Erik Mickoeit et al. que teóricamente no destacar respecto a las demás por su COP pero es la única que consigue tener un VHC más alto que el propano.

Analizando más en detalle las tres mezclas con más COP mediante figuras ternarias donde se expresan diferentes magnitudes en función de la composición de la mezcla, se puede observar que en las dos primeras, el punto de máximo rendimiento que ha aparecido en el resultado es una región muy pequeña y muy sensible a variaciones de la composición (\autoref{fig:COP_alta_n1} y \autoref{fig:COP_alta_n2}), esto no es deseable ya que si hay fugas en la instalación y la composición varía no se quiere que el comportamiento se vea afectado drásticamente. Además, estas regiones de máximo rendimiento son regionnes límite, esto significa que si varía la concentración un poco podríamos estar en puntos de trabajo que no cumplen con las restricciones que se han establecido previamente, en este caso, la variable que no cumpliría es el glide en el evaporador, este es muy sensible a cambios en la composición y podría superar los 10ºC (\autoref{fig:glide_o_alta_n1}, \autoref{fig:glide_o_alta_n2}).

\begin{figure}[ht]
    \centering
    \begin{minipage}[t]{0.48\textwidth}
        \centering
        \includegraphics[width=\textwidth]{imagenes/res_teo/diag_ternarios/alta/PROPYLENE_DME_BUTANE/COP.png}
        \caption{\centering COP en función de la composición}
        \label{fig:COP_alta_n1}
    \end{minipage}
    \hfill
    \begin{minipage}[t]{0.48\textwidth}
        \centering
        \includegraphics[width=\textwidth]{imagenes/res_teo/diag_ternarios/alta/PROPYLENE_DME_BUTANE/glide 0.png}
        \caption{\centering Glide en el evaporador en función de la composición}
        \label{fig:glide_o_alta_n1}
    \end{minipage}
\end{figure}

\begin{figure}[ht]
    \centering
    \begin{minipage}[t]{0.48\textwidth}
        \centering
        \includegraphics[width=\textwidth]{imagenes/res_teo/diag_ternarios/alta/PROPANE_DME_BUTANE/COP.png}
        \caption{\centering COP en función de la ocmposición}
        \label{fig:COP_alta_n2}
    \end{minipage}
    \hfill
    \begin{minipage}[t]{0.48\textwidth}
        \centering
        \includegraphics[width=\textwidth]{imagenes/res_teo/diag_ternarios/alta/PROPANE_DME_BUTANE/glide 0.png}
        \caption{\centering Glide en el evaporador en función de la composición}
        \label{fig:glide_o_alta_n2}
    \end{minipage}
\end{figure}

\FloatBarrier

Siguiendo con la siguiente mezcla (R1270, RE170, R600a) en esta se puede observar un comportamiento diferente, las regiones de máximo rendimiento ocupan una zona más amplia y esta no es una frontera con lo que hay margen con posibles variaciones de la composición (\autoref{fig:COP_alta_n3}). Esta zona de alto COP se puede explicar por la subida de presiones en el evaporador que conllevan una reducción en la tasa de compresión y por ende, una reducción de trabajo del compresor (\autoref{fig:presion_o_alta_n3}).

\begin{figure}[ht]
    \centering
    \begin{minipage}[t]{0.48\textwidth}
        \centering
        \includegraphics[width=\textwidth]{imagenes/res_teo/diag_ternarios/alta/PROPYLENE_DME_ISOBUTANE/COP.png}
        \caption{\centering COP en función de la composición}
        \label{fig:COP_alta_n3}
    \end{minipage}
    \hfill
    \begin{minipage}[t]{0.48\textwidth}
        \centering
        \includegraphics[width=\textwidth]{imagenes/res_teo/diag_ternarios/alta/PROPYLENE_DME_ISOBUTANE/Presion 0.png}
        \caption{\centering Presión de evaporación en función de la composición}
        \label{fig:presion_o_alta_n3}
    \end{minipage}
\end{figure}

\FloatBarrier

En la figura \autoref{fig:PH_3_mejores_alta} se pueden observar la comparación de las tres mezclas ternarias respecto al propano.

\begin{figure}[ht]
    \centering
    \includegraphics[width=\textwidth]{imagenes/res_teo/diagramas_PH/alta/3_mejores_con.png}
    \caption{Diagramas P-H de las mejores mezclas}
    \label{fig:PH_3_mejores_alta}
\end{figure}

\FloatBarrier

\hl{ME PAREIX QUE SE VEU UN POC MALAMENT, HI HA MASSA LINIES, ¿POSO NOMÉS EL DIAGRAMA DE LES MESCLES SENSE ELS PUNTS DEL CICLE?}

Continuando con el análisis siguen las mezclas que se componen por su mayoría de Dimetíl Éter (RE170), esto se debe a que, entre otros motivos, el calor latente del Dimetil Éter es muy grande. Este fonómeno se puede observar en la \autoref{fig:PH_fluidos_no_volatiles} donde se comparan los diagramas P-H de los diferentes fluidos con una volatilidad similar al propano. \hl{¿FA FALTA QUE POSE PER QUÈ HA DE SER AMB VOLATILITAT SIMILAR AL PROPÀ?}

 % (PER SI FA FALTA) Solo se comparan los fluidos con una valatilidad similar ya que si se comparara fluidos de diferente volatilidad las presiones en el ciclo variarían demasiado pudiendo llegar a ser demasiado altas o demasiado bajas, por lo tanto, estos fluidos siempre estarán contenidos a una fracción pequeña del ciclo.

\begin{figure}[ht]
    \centering
    \includegraphics[width=\textwidth]{imagenes/res_teo/diagramas_PH/fluidos_puros_no_volatiles.png}
    \caption{Comparación de fluidos con volatilidad similar al R290}
    \label{fig:PH_fluidos_no_volatiles}
\end{figure}

\FloatBarrier

Como todos las mezclas que combinan con el RE170 tienen un incremento del COP respecto al R290 similar rondando los 10.8\%, se pasará a mirar el segundo criterio más importante, la VHC.

En la \autoref{tab:res_teo_vhc} se pueden observar los fluidos ordenados por incremento en la capacidad calorífica volumétrica.

\begin{table}
    \centering
    
\begin{tabular}{|cc|cc|c|c|cc|}
\hline
\multicolumn{2}{|c|}{\multirow{2}{*}{\textbf{Fluidos}}} & \multicolumn{2}{c|}{\multirow{2}{*}{\textbf{Proporciones}}} & \multirow{2}{*}{\textbf{$\Delta$VHC}} & \multirow{2}{*}{\textbf{$\Delta$COP}} & \multicolumn{2}{c|}{\textbf{Glide (ºC)}} \\ \cline{7-8} 
\multicolumn{2}{|c|}{} & \multicolumn{2}{c|}{} &  &  & \multicolumn{1}{c|}{\textbf{cond}} & \textbf{evap} \\ \hline
\multicolumn{1}{|c|}{\textbf{RE170}} & \textbf{R1270} & \multicolumn{1}{c|}{\textbf{93\%}} & \textbf{7\%} & \textbf{-23.1\%} & \textbf{+10.9\%} & \multicolumn{1}{c|}{\textbf{1.1}} & \textbf{0.8} \\ \hline
\multicolumn{1}{|c|}{RE170} & R290 & \multicolumn{1}{c|}{95\%} & 5\% & -24.1\% & +10.8\% & \multicolumn{1}{c|}{1.0} & 0.7 \\ \hline
\multicolumn{1}{|c|}{RE170} & R600 & \multicolumn{1}{c|}{99\%} & 1\% & -25.1\% & +10.2\% & \multicolumn{1}{c|}{0.1} & 0 \\ \hline
\multicolumn{1}{|c|}{RE170} & R170 & \multicolumn{1}{c|}{98\%} & 2\% & -25.1\% & +10.8\% & \multicolumn{1}{c|}{1.4} & 0.5 \\ \hline
\multicolumn{1}{|c|}{RE170} & R1150 & \multicolumn{1}{c|}{99\%} & 1\% & -25.9\% & +10.6\% & \multicolumn{1}{c|}{1.5} & 0.4 \\ \hline
\multicolumn{1}{|c|}{RE170} & R744 & \multicolumn{1}{c|}{99\%} & 1\% & -26.0\% & +10.9\% & \multicolumn{1}{c|}{1.2} & 0.3 \\ \hline
\multicolumn{2}{|c|}{RE170} & \multicolumn{2}{c|}{100\%} & -27.8\% & +9.7\% & \multicolumn{1}{c|}{0} & 0 \\ \hline
\end{tabular}

    \caption{Mezclas binarias ordenadas por VHC}
    \label{tab:res_teo_vhc}
\end{table}

\newpage
\end{document}